\scanpage{158}
\mspage{74}
% archive sheet number: 157

Condition a) : $\mathrm{Filt}^\pi$ équivalente à une petite catégorie : c'est immédiat, puisque \unsure{par} (4.3.6.2) $\mathrm{Filt}^\pi$ est formé des $\underset{I}{\unsure{\underline{\lim}|\underrightarrow{\lim}}}$ \inserted{dans $M$} d'objets de $S$ (petite catégorie), avec $I$ telles que $\mathrm{card}\,\mathrm{Fl}\,I\leq\pi$.

Condition b) : $\mathrm{Filt}^\pi$\textsuperscript{*} stable par $\underset{I}{\underrightarrow{\lim}}$, $\mathrm{card}\,\mathrm{Fl}\,I\leq\pi$. Soit donc

% slanted
\margin{\unsure{prendre|prend} $I$ filtrant \unsure{gr.|et gr.} devant $\pi_0$ \unsure{je ne} le \unsure{prouve} que dans ce cas, hélas !}

\[
(x_i)_{i\in I},\ \deleted{\unsure{\text{ou}|\text{un}}}\qquad I\mathrel{\mathop{\longrightarrow}\limits^{\unsure{\varphi|\rho}}_{i\,\unsure{\mapsto|\vdash}\,x_i}}M
\]

un syst. inductif dans $M$, avec des $x_i$ dans $\mathrm{Filt}^\pi$. Considérons le produit fibré

\diagram{\[
\begin{tikzcd}[ampersand replacement=\&,column sep=large,row sep=large]
I\arrow[d,dash,"\varphi"'] \& S_{/I}=J\arrow[l,"\text{cofibré}"']\arrow[d]\arrow[dr,dashed,"\psi"] \& \\
M \& S_{/M}\arrow[l,"\text{but}"',"\text{Cat-cofibrée}"]\arrow[r,"\text{source}"] \& S\unsure{\lhook\joinrel-|\subset|\hookrightarrow}M
\end{tikzcd},
\tag*{(4.3.6.4)}
\]}

où $J$ est Cat-cofibré sur $I$, sa fibre en $i\in\mathrm{Ob}\,I$ étant isom. à $S_{/x_i}$, donc $\mathrm{Card}\,\mathrm{Fl}\,\unsure{I_i|J_i}\leq\pi$. Comme par hypothèse $\mathrm{card}\,\mathrm{Fl}\,I\leq\pi$, on en conclut que $\mathrm{card}\,\mathrm{Fl}\,J\leq\pi^3=\pi$,

[NB L'ens. des flèches \inserted{$v:j_0\longrightarrow j_1$} de $J$ au dessus d'une flèche $u:i_0\longrightarrow i_1$ de $I$ est ``contenu'' dans l'ens. produit $\mathrm{Ob}\,I_{i_0}\times\mathrm{Fl}\,I_{i_1}$, par l'application $v\mathrel{\unsure{\mapsto|\vdash}}(\text{source }v,v':u_*(\unsure{i_0|j_0})\longrightarrow j_1)$ donc de cardinal majoré par $\pi^2$.] \deleted{\unsure{Soit} \unsure{$x=\lim$}} On

